% Generated draft: requires mathematical and source review.
\documentclass[11pt]{article}
\usepackage[margin=25mm]{geometry}
\usepackage{fontspec}
\setmainfont{Latin Modern Roman}
\usepackage{amsmath,amssymb,mathtools,unicode-math}
\setmathfont{Latin Modern Math}
\usepackage{enumitem,needspace,xurl,hyperref}
\setlength{\parindent}{0pt}
\setlength{\parskip}{5pt}
\setlength{\emergencystretch}{3em}
\newcommand{\N}{\mathbb{N}}
\newcommand{\Z}{\mathbb{Z}}
\newcommand{\Q}{\mathbb{Q}}
\newcommand{\R}{\mathbb{R}}
\newcommand{\C}{\mathbb{C}}
\newcommand{\F}{\mathbb{F}}
\newcommand{\A}{\mathbb{A}}
\newcommand{\PP}{\mathbb{P}}
\newcommand{\E}{\mathbb{E}}
\newcommand{\eps}{\varepsilon}
\newcommand{\dd}{\,\mathrm{d}}
\begin{document}
\section*{UnsolvedMath ID: 700378}
\textbf{Source problem code:} catalogue entry
\subsection{Definitions and mathematical statement}
Let \(n\) be an arbitrary integer with \(2\leq n<\infty\), and let \(F_n\) be the free profinite group of rank \(n\). Write \(F_{\aleph_0}\) for the free profinite group of countably infinite rank intended by the source.

Under the interpretation that normal subgroups are closed and isomorphisms are isomorphisms of topological groups, the question is whether
\[
\forall N\trianglelefteq F_{\aleph_0}\text{ closed},\quad
\exists M\trianglelefteq F_n\text{ closed}\quad
N\cong_{\mathrm{top}} M.
\]
Here each subgroup carries the subspace topology, and \(N\cong_{\mathrm{top}}M\) means that there is a group isomorphism \(N\to M\) that is continuous with continuous inverse.

The supplied statement does not explicitly require the normal subgroups to be closed or specify that the isomorphism is topological. Consequently, the displayed formulation is conditional on those conventions. Allowing arbitrary abstract normal subgroups or requiring only abstract group isomorphisms gives different formulations that the supplied record does not distinguish. The interpretation of ``countable rank'' as countably infinite rank, and the precise free profinite basis convention at that rank, also require confirmation from the source.
\subsection{Short English statement}
For every finite rank \(n>1\), must every normal subgroup of the free profinite group of countable rank occur, up to isomorphism, as a normal subgroup of \(F_n\)? The formulation above uses closed normal subgroups and topological isomorphisms; these conventions require source confirmation.
\subsection{Sources}
\begin{itemize}
\item Local UnsolvedMath registry, supplied complete source record: permanent ID \(700378\), exact source problem code \texttt{catalogue entry}, title \emph{Kourovka Notebook Problem 6.32}. Upstream revision: \texttt{d371ee917ac9647fbb06cd6fc529deec6f3dda82}.
\item Supplied attribution: O. V. Mel'nikov, proposed in 1978; \emph{The Kourovka Notebook (Unsolved Problems in Group Theory)}, 6th issue (1978), Problem 6.32. The record also identifies the 21st issue edition (2026), \texttt{arXiv:1401.0300v48}. The supplied collection description names E. I. Khukhro and V. D. Mazurov as editors.
\item The local record labels the problem \emph{open} and gives an update timestamp of \texttt{2026-10-16T14:00:00Z}. This is supplied registry metadata, not an independently verified mathematical or literature assessment. The cited source has not been independently consulted for this draft.
\end{itemize}
\end{document}
